% The 200 word abstract, inside the window the journal states.
% Swap it into paper/main.tex if the 251 word version is questioned at production.
\begin{abstract}
A denoiser for optical coherence tomography minimises a pixel error, so it can raise fidelity while
lowering the layer contrast and boundary sharpness a clinician reads. We repair those properties
afterwards, with the denoiser frozen. Four things are new. Six clinical properties in closed form read
the noisy input and the denoised output alone, each returning a score and a map of where it fails, so
no clean image is needed at test time. Five fuzzy rules turn those maps into a per-pixel allocation
whose response to a shift in evidence is bounded by an exact rather than a loose constant. A bounded
proposal is admitted in proportion to how many of three constraints hold, we prove when contrast
cannot fall, and every change traces to a named property at a named place. One set of rules and
weights serves four denoisers, only a small head refitted. All six measures rise on two of them and
five of six on the others, contrast by 3.5 to 11.3 percent, speckle suppression by 23.7 to 57.1 and
signal to noise by 8.0 to 23.0, at 0.17 to 0.98 dB of fidelity. Carried unchanged to two other
datasets, 40 of 48 measures improve.
\end{abstract}
